\documentclass[11pt]{article}
\usepackage{amsmath}
\usepackage{amssymb}
\usepackage{geometry}
\geometry{margin=1in}

\title{Plan 01-03: J3 Potential in Spherical Coordinates}
\author{J2+J3+Drag Secular Perturbation Theory}
\date{2026-04-02 — Phase 1}

\begin{document}
\maketitle

\section{Objective}
Express Earth's J3 gravitational perturbation potential in spherical coordinates.

\section{J3 Gravitational Potential}

The third zonal harmonic J3 contributes:

\begin{equation}
U_{J3}(r, \phi) = -\frac{\mu}{r}\left[J_3\left(\frac{R_E}{r}\right)^3 P_3(\sin \phi)\right]
\end{equation}

where:
\begin{itemize}
    \item $J_3 = -2.5327 \times 10^{-6}$ (negative, Northern hemisphere bulge)
    \item $P_3(x) = (5x^3 - 3x)/2$ (Legendre polynomial, degree 3)
\end{itemize}

\section{Disturbing Potential}

\begin{equation}
\boxed{R_{J3}(r, \phi) = \frac{\mu J_3 R_E^3}{r^4} P_3(\sin \phi)}
\label{eq:rj3}
\end{equation}

Expanded:
\begin{equation}
R_{J3}(r, \phi) = \frac{\mu J_3 R_E^3}{2r^4}\left(5\sin^3 \phi - 3\sin \phi\right)
\end{equation}

\section{Key Differences from J2}

\begin{enumerate}
    \item \textbf{Odd degree:} $r^{-4}$ vs. $r^{-3}$ for J2
    \item \textbf{Odd function:} $P_3(-x) = -P_3(x)$ creates North-South asymmetry
    \item \textbf{Magnitude:} $|J_3/J_2| \approx 2.3 \times 10^{-3}$ (J3 is ~0.23\% of J2)
\end{enumerate}

\section{Dimensional Analysis}

\begin{equation}
[R_{J3}] = \frac{[\mu][R_E^3]}{[r^4]} = \frac{\text{m}^3/\text{s}^2 \cdot \text{m}^3}{\text{m}^4} = \text{m}^2/\text{s}^2 \quad \checkmark
\end{equation}

\section{Success Criteria}
\begin{itemize}
    \item[$\checkmark$] J3 potential with $P_3$ Legendre polynomial
    \item[$\checkmark$] Odd-degree structure noted
    \item[$\checkmark$] Dimensional analysis passed
\end{itemize}

\end{document}
